\documentclass[12pt]{article}
\usepackage{graphicx} % Use this package to include images
\usepackage{amsmath} % A library of many standard math expressions
\usepackage[margin=1in]{geometry}% Sets 1in margins.
\usepackage{fancyhdr} % Creates headers and footers
\usepackage{enumerate}  %These two package give custom labels to a list
\usepackage{enumitem}
\usepackage{amssymb}
\usepackage{amsthm} % for theorem style environments for question and answers for readablity
\usepackage{titlesec}
\usepackage[backend=biber, style=numeric]{biblatex}
\addbibresource{biblio.bib}
% Customize section font size
\titleformat*{\section}{\fontsize{14pt}{16pt}\bfseries\selectfont}
% Customize subsection font size
\titleformat*{\subsection}{\fontsize{12pt}{14pt}\bfseries\selectfont}
\titlespacing*{\section}{0pt}{1ex}{0.5ex}
\titlespacing*{\subsection}{0pt}{0.8ex}{0.4ex}
\newtheorem*{thm}{Theorem}
\theoremstyle{definition}
\newtheorem*{q}{Question}
\newtheorem*{ans}{Answer}
\title{Outline Draft}
\author{Loveless, Nix, Stokes, Abernathy's}
\date{\today}
\begin{document}
\maketitle
\vspace{-2em}  % Remove space after title block
\section{Introduction}
Alzheimer's Disease (AD) is characterized as a complex neurodegenerative disease affecting approximately more than 55 million people globally \cite{maji2025mathematical}. 
Symptoms are caused by the degeneration of the central nervous system, and are portrayed by progressive memory loss, difficulty with language, drastic mood swings, among other effects.
It is a disease that causes an extensive burden to those suffering from the disease through disorientation, to relatives in the form of mental taxation, as well as economically; 
projecting an estimated global economic burden of $16.9$ trillion dollars by 2050 \cite{nandi2022global}.
These factors are but a few reasons as to the importance of research in this area, where many scientists have worked to understand the underlying mechanisms of AD.
However, much is still unknown in how AD is caused, as well as what biological processes occur in the brain, making the disease difficult to come to a decisive conclusion.

Many studies have been conducted in regard to understanding Alzheimer's Disease. 
With this, there are a few main hypotheses that are generally considered to be factors in the causation and progression of AD \cite{zhang2024recent}.
The so-called \textit{Amyloid Cascade Hypothesis} is among the most extensively studied in the literature, where it is proposed that a misfolding of $\beta$-amyloid (A $\beta$) protein form
senile plaques. These plaques are hypothezied to be the main driving force in the intracellular deposition of misfolded tau proteins in neurofibrillary tangles (NFTs) and memory loss in the 
progression of Alzheimer's Disease \cite{hardy1992alzheimer}. Growing evidence however, shows that this hypothesis is suggested to be but a small factor in the larger picture of the underlying 
operative mechanisms of the disease \cite{kepp2023amyloid}. This has led researchers to branch into other pathways concerning AD, looking at the interactions between A $\beta$ and tau in an attempt
to further understand the pathogensis of these two factors \cite{zhang2021interaction}.




















\section{Equations}
Potential systems consist of various mechanisms, but the main differences come from whether or not we contain a Holling-Type III term in our $A_ \beta$ compartment. 
The other difference comes from whether or not we include a logistic growth term in each compartment \cite{petrella2019computational}, \cite{pal2025nonequilibrium}.
The Equations including the Holling-Type III term is given by 
\begin{eqnarray}
    \frac{dA _{\beta}}{dt}&=&\lambda _{A _{\beta}}A _{\beta}(K _{A _{\beta}}-A _{\beta})+\lambda _{A _{\beta Ca}}\frac{Ca ^{2}}{Ca ^{2}+\sigma ^{2}}-(1-\alpha _{A _{\beta}})\delta _{A _{\beta}}A _{\beta}\\
    \frac{dCa}{dt}&=& \lambda _{Ca}Ca (K _{Ca}-Ca)+\lambda _{CaA _{\beta}}-(1-\alpha _{Ca})\delta _{Ca} Ca\\
    \frac{d \tau}{dt}&=&\lambda _{\tau}\tau (K _{\tau}-\tau)+\lambda _{\tau A _{\beta}}A _{\beta}+\lambda _{\tau Ca}-(1-\alpha _{\tau})\delta _{\tau}\tau\\
    \frac{dN}{dt}&=&\lambda _{N}N (K _{N}-N) + \lambda _{N\tau} \tau\\
    \frac{dC}{dt}&=&\lambda _{C}C (K _{C}-C)+\lambda _{CN}NR+\lambda _{C\epsilon}\epsilon.
\end{eqnarray}

A simpler model can be defined, excluding the Holling-Type III term within $dA _{\beta}/dt$.
\begin{eqnarray}
    \frac{dA _{\beta}}{dt}&=&a _{1}-k _{1}A _{\beta}+\lambda _{A _{\beta}Ca}-(1-\alpha _{A _{\beta}})\delta _{A _{\beta}}A _{\beta}\\
    \frac{dCa}{dt}&=&a _{2} -k _{2}Ca+\lambda _{CaA _{\beta}}-(1-\alpha _{Ca})\delta _{Ca} Ca\\
\frac{d \tau}{dt}&=&a _{3}-k _{3}\tau+\lambda _{\tau A _{\beta}}A _{\beta}+\lambda _{\tau Ca}-\underbrace{(1-\alpha _{\tau})\delta _{\tau}\tau}_{\text{off the dome}}\\
    \frac{dN}{dt}&=&a _{4}-k _{4}N+ \lambda _{N\tau} \tau\\
    \frac{dC}{dt}&=&\underbrace{a _{5}-k _{5}C}_{\text{Pal}}+\underbrace{\lambda _{CN}NR+\lambda _{C\epsilon}\epsilon}_{\text{Petrella}}.
\end{eqnarray}

In each system, $a _{n},k _{n}$ is the growth and wash out, $\lambda _{ab}$ is the growth rate for $a$ with respect to $b$, $ab \neq 0$. 
$\alpha _{n}$ is the inefficiency of some drug for $n$, $0<\alpha _{n}<1$.
$\delta _{n}$ is the strength of some drug, $0<\delta _{n}<1$.

The proposed \textbf{MODEL TO END ALL MODELS}:
\begin{eqnarray}
    \frac{dA _{\beta}}{dt}&=&a _{1}-r _{1}A _{\beta}+a _{2}\frac{Ca ^{2}}{Ca ^{2}+\sigma ^{2}}\\
    \frac{dCa}{dt}&=& b _{1}-r _{2}Ca+b _{2}A _{\beta}\\
    \frac{d \tau}{dt}&=&c _{1}-r _{3}\tau+c _{2}A _{\beta}+c _{3}Ca\\
    \frac{dN}{dt}&=&d _{1}-k _{4}N+d _{2}\tau\\
    \frac{dC}{dt}&=&e _{1}-k _{5}C+e _{2}NR.
\end{eqnarray}

where $r _{n}= k _{n}+(1-\alpha _{n})\delta _{n}$. 
\subsection{Week 3 Equations Currently in work}
\begin{eqnarray}
    \frac{dA _{\beta}}{dt}&=&a _{1}-r _{1}A _{\beta}+a _{2}\frac{Ca ^{2}}{Ca ^{2}+\sigma ^{2}}\\
    \frac{dCa}{dt}&=& b _{1}-r _{2}Ca+b _{2}A _{\beta}\\
    \frac{d \tau}{dt}&=&c _{1}-r _{3}\tau+c _{2}A _{\beta}+c _{3}Ca\\
    \frac{dN}{dt} &=& d _{1}\tau \left(1 - \frac{N}{K _{N}}\right)\\
    \frac{dC}{dt} &=& (e _{1}N+e _{2}\tau)\left(1 - \frac{C}{K _{C}}\right)
\end{eqnarray}
and
\begin{eqnarray}
    \frac{dA _{\beta}}{dt}&=&a _{1}-r _{1}A _{\beta}+a _{2}\frac{Ca }{Ca+\sigma}\\
    \frac{dCa}{dt}&=& b _{1}-r _{2}Ca+b _{2}A _{\beta}\\
    \frac{d \tau}{dt}&=&c _{1}-r _{3}\tau+c _{2}A _{\beta}+c _{3}Ca\\
    \frac{dN}{dt}&=&d _{1}-k _{4}N+d _{2}\tau\\
    \frac{dC}{dt}&=&e _{1}-k _{5}C+e _{2}NR + e _{3}\tau.
\end{eqnarray}


\subsection{Mass Action Terms Model}
\begin{eqnarray}
    \frac{dA _{\beta}}{dt}&=&a _{1}-c _{1}A _{\beta}-k _{1}A _{\beta}\tau - \alpha _{1}A _{\beta}\\
    \frac{dCa}{dt}&=& b _{1}+ k _{2} A _{\beta}\tau - c _{2}\tau - \alpha _{2} \tau\\
    \frac{d \tau}{dt}&=&k _{3}A _{\beta}+k _{4}\tau - c _{3}C _{a}- \alpha _{3}C _{a}\\
    \frac{dN}{dt}&=&rN \left(1 - \frac{N}{K}\right)- k _{5}C _{a}N- k _{6}\tau N.
\end{eqnarray}

\section{Parameters}
\begin{table}
    \begin{tabular}{l|l|l|l}
        Var & Value & Unit & Source \\
        \hline
        $A _{\beta}$ & $\sim 0.01$ & $C$& \cite{petrella2019computational}\\
        $Ca$ & $0$& &\\
        $\tau$ & $0$& $C$ & \cite{petrella2019computational} \\
        $N$ & $0$ & &\\
        $C$ & $0$ & &\\
        $a _{1}$ & $0.25$ & $c/yr$& \cite{pal2025nonequilibrium} \\
        $a _{2}$ & $2.89$ & $c/yr$& \cite{pal2025nonequilibrium}\\
        $b _{1}$ & $11$ & $c/yr$ & \cite{pal2025nonequilibrium}\\
        $b _{2}$ & $1000$ & $yr ^{-1}$ & \cite{pal2025nonequilibrium}\\
        $c _{1}$ & & &\\
        $c _{2}$ & $0.15$ & $yr ^{-1}$& \cite{hao2022optimal}\\
        $c _{3}$ & & &\\
        $d _{1}$ & & &\\
        $d _{2}$ & & &\\
        $e _{1}$ & $1.67$ & $yr _{-1}$ & \cite{hao2022optimal}\\
        $e _{2}$ & $3.83$ &$yr _{1}$ &\cite{hao2022optimal}\\
        $k _{1}$ & $0.35$ & $yr ^{-1}$ & \cite{pal2025nonequilibrium}\\
        $k _{2}$ & $500$ & $yr ^{-1}$ & \cite{pal2025nonequilibrium}\\
        $k _{3}$ & & &\\
        $k _{N}$ & $1.00$ & $C$ &\\
        $k _{C}$ & $169.48$ & $C$ &\\
        $\alpha _{A _{\beta}}$ & & &\\
        $\alpha _{Ca}$ & & &\\
        $\alpha _{\tau}$ & & &\\
        $\delta _{A _{\beta}}$ & & &\\
        $\delta _{Ca}$ & & &\\
        $\delta _{\tau}$ & & &\\
        $\sigma$ & $0.00027$& $C$ & \cite{zhang2020mathematical}
    \end{tabular}
\end{table}




\printbibliography

\end{document}
